\section{Resultados}

\subsection{Reportes de HDFS}
Durante la ejecuci\'on del laboratorio, se obtuvieron diversos reportes provenientes del sistema de archivos distribuido HDFS. Los comandos utilizados permitieron verificar tanto el estado general del sistema como la presencia y distribuci\'on de los datos almacenados. Se observ\'o que el NameNode registr\'o correctamente la estructura del almacenamiento distribuido, mientras que los DataNodes reportaron disponibilidad operativa y espacio asignado.

Asimismo, la inserci\'on de archivos de prueba evidenci\'o que el mecanismo de replicaci\'on funcion\'o conforme a los par\'ametros predefinidos. HDFS mantuvo copias redundantes que garantizaron tolerancia a fallos, lo que se comprob\'o mediante consultas posteriores y verificaciones de integridad.

\subsection{Salidas de MapReduce}
Las ejecuciones de los trabajos MapReduce generaron salidas consistentes con los algoritmos implementados. Entre los resultados observados se incluyeron:

\begin{itemize}
    \item Registros de mapeo correctamente distribuidos entre nodos disponibles.
    \item Reducci\'on exitosa de claves y valores, consolidando los resultados parciales generados durante la fase de mapeo.
    \item Producci\'on de archivos de salida ubicados en los directorios configurados dentro de HDFS.
\end{itemize}

Los logs de MapReduce indicaron que las tareas se ejecutaron sin interrupciones y con asignaci\'on equilibrada de recursos, confirmando la operatividad del JobTracker y TaskTracker (o sus equivalentes en YARN si la configuraci\'on era posterior).

\subsection{Rendimiento observado}
El rendimiento general del cl\'uster construido sobre Docker mostr\'o estabilidad durante el procesamiento de datos. A pesar de tratarse de un entorno virtualizado, los tiempos de ejecuci\'on se mantuvieron dentro de los valores esperados para un laboratorio controlado.

Se presentaron los siguientes hallazgos clave:

\begin{itemize}
    \item El tiempo total de ejecuci\'on de MapReduce fue consistente a trav\'es de diferentes pruebas.
    \item No se registraron fallos de nodos ni desconexiones inesperadas.
    \item El consumo de CPU y memoria se mantuvo dentro de l\'imites aceptables para contenedores.
\end{itemize}

En conjunto, los resultados indican que el entorno Docker proporcion\'o una plataforma adecuada para simular un cl\'uster Hadoop funcional.
